\section{The classification of second-order elementary automata}\label{sec:classification}

By Lemma~\ref{lem:classes}, $\Aut(X_{F_f})$ and every statement below about block maps of bounded radius depend only on the class of $f$ under reflection and the cellwise complement. This section states what each status means, lists the classes with the results that settle them, proves Theorem~\ref{thm:main}, and collects what the computations say about the open classes.

\subsection{What the entries mean}\label{cl:sec:status}

\begin{definition}\label{cl:def:status}
A class has one of the following five statuses.
\begin{enumerate}
\item[(G)] \emph{Every finite group.} There is an injective homomorphism $\bigoplus_{m\ge1}S_m\to\Aut(X_{F_f})$. Then $\Aut(X_{F_f})$ contains every finite group and is not virtually $\Z^n$ for any $n$. This is a lower bound. Only for the classes of $255$ and $95$ is the whole group described: for $255$ as $\Z/4\times\mathcal G_4$ \eqref{imp:affine}, and for $95$ in terms of the group $\mathcal N$ \eqref{imp:95}.
\item[(Z)] \emph{$\Z^2$.} $\Aut(X_{F_f})=\langle\sh,F_f\rangle$, and $(a,b)\mapsto\sh^aF_f^b$ is an isomorphism $\Z^2\to\langle\sh,F_f\rangle$.
\item[(S)] \emph{$\Z^2\times\Z/2$.} $\Aut(X_{F_f})=\langle\sh,F_f\rangle\times\langle\kappa\rangle\cong\Z^2\times\Z/2$, where $\kappa$ is the cellwise complement.
\item[(A)] \emph{Affine, free abelian.} $f$ is affine, and $\Aut(X_{F_f})$ is a free abelian group of finite rank that contains $\langle\sh,F_f\rangle$ as a proper subgroup.
\item[(O)] \emph{Open.} None of the above is known.
\end{enumerate}
\end{definition}

The statuses (G), (Z), (S) and (A) exclude one another: a group of status (Z), (S) or (A) is virtually $\Z^n$, a group of status (A) is torsion-free and larger than $\langle\sh,F_f\rangle$. The status (O) is a statement about what is known, not about the automata; its mathematical content is the census of Theorem~\ref{cl:thm:census}, which finds no automorphism outside $\langle\sh,F_f\rangle$ of radius at most $3$. The census says nothing about larger radii.

\begin{remark}\label{cl:rem:affine}
For an affine rule $f(a,b,c)=\alpha a\oplus\beta b\oplus\gamma c\oplus\varepsilon$, let $\mathcal R=\F_2[s,s^{-1}]$ act on $\F_2^{\Z}$ through $s=\sh$. Then $F_f(c)=Mc+k_0$, where $M=\begin{psmallmatrix}0&1\\1&p\end{psmallmatrix}$ with $p=\alpha s^{-1}+\beta+\gamma s$ acts on columns $(x,y)^{\mathsf T}$ and $k_0=(0,\varepsilon)^{\Z}$. An affine block map $c\mapsto Nc+k$, with $N$ linear and $k$ a constant configuration, commutes with $F_f$ if and only if $NM=MN$ and $(N+I)k_0=(M+I)k$, and the matrices commuting with $M$ are the $aI+bM$ with $a,b\in\mathcal R$. There are nine affine classes. For those of $15$, $60$, $105$ and $153$ every automorphism is of this form and the groups are determined \eqref{imp:affine}; those of $0$, $51$, $90$, $165$ and $255$ have status (G). For $255$, $f\equiv1$ and $F_{255}(x,y)=(y,\neg x)$ is the cellwise permutation $0\mapsto1\mapsto3\mapsto2\mapsto0$, so $F_{255}^2=\kappa$, $F_{255}^4=\mathrm{id}$ and $\langle\sh,F_{255}\rangle\cong\Z\times\Z/4$.
\end{remark}

\subsection{The group generated by the shift and the automaton}\label{cl:sec:Z2}

For an elementary rule $f$ let $\Phi_f\subseteq\{0,1\}^{\Z}$ be the set of sequences $\varphi$ with
\[
\varphi(k-1)\oplus\varphi(k+1)=f\bigl(\varphi(k-1),\varphi(k),\varphi(k+1)\bigr)\qquad\text{for all }k\in\Z.
\]

\begin{lemma}[{\cite[Lemma~2.3]{PaperA}}]\label{cl:lem:Z2}
The maps $\varphi\mapsto c$, $c_i=(\varphi(i+1),\varphi(i))$, and $\varphi\mapsto d$, $d_i=(\varphi(i-1),\varphi(i))$, are bijections from $\Phi_f$ onto the fixed points of $\sh F_f$ and of $\sh^{-1}F_f$ respectively.
\end{lemma}

The set $\Phi_f$ is the shift of finite type on $\{0,1\}$ given by the allowed triples $(u,v,w)$ with $u\oplus w=f(u,v,w)$. It has a point that is not periodic if and only if its graph on pairs has a strongly connected component that is not a single cycle, or a path between two distinct cycles; this is a finite check, and its outcome is recorded in Tables~\ref{cl:tab:decided} and~\ref{cl:tab:open}. For $110$ the allowed triples are $000,001,011,110,111$, and $\Phi_{110}$ consists of $0^{\Z}$, $1^{\Z}$ and the steps $0^{(-\infty,p]}1^{(p,\infty)}$. We call a point of $\mathrm{Fix}(\sh^aF_f^b)$ that is not spatially periodic but agrees with spatially periodic configurations on a left and on a right half-line a \emph{kink}; the fixed points of $\sh^{\pm1}F_{110}$ that correspond to the steps are the kinks on which the proof of \eqref{imp:110} is built \cite{PaperA}. If $\Phi_f$ contains a sequence that is not periodic, then the corresponding $c$ and $d$ are not periodic under $\sh$, and $\sh^aF_f^b=\mathrm{id}$ gives $c=\sh^{a-b}c$ and $d=\sh^{a+b}d$, so $a=b=0$ and $\langle\sh,F_f\rangle\cong\Z^2$. The following theorem gives this without a hypothesis on $\Phi_f$.

A block map $g$ \emph{depends on} a bit of the cell $c_i$ (for the output bit $y(g(c)_0)$, say) if two configurations that differ only in that bit have different outputs. For a nonconstant block map $g$ let $r^+(g)$ be the largest and $-r^-(g)$ the smallest $i$ such that $g(c)_0$ depends on $c_i$.

\begin{theorem}[{\cite[Theorem~2.4]{PaperC}}]\label{cl:thm:A}
Let $f$ be an elementary rule, $F=F_f$ and $n\ge1$.
\begin{enumerate}
\item If $f$ depends on its third argument, then $y(F^n(c)_0)$ depends on $y(c_n)$, and $x(F^{-n}(c)_0)$ depends on $x(c_n)$. So $r^+(F^n)=r^+(F^{-n})=n$.
\item If $f$ depends on its first argument, then $y(F^n(c)_0)$ depends on $y(c_{-n})$, and $x(F^{-n}(c)_0)$ depends on $x(c_{-n})$. So $r^-(F^n)=r^-(F^{-n})=n$.
\item The map $(a,b)\mapsto\sh^aF^b$ is an isomorphism $\Z^2\to\langle\sh,F\rangle$ if and only if $f$ is not a function of its middle argument alone.
\end{enumerate}
\end{theorem}

The rules that are functions of their middle argument alone are $0$, $51$, $204$ and $255$; for them $F$ acts by the same permutation of $\A$ in every cell, of order $2$, $3$, $3$ and $4$ respectively. So $F_f$ has infinite order and $\langle\sh,F_f\rangle\cong\Z^2$ in every class except the three classes $\{0\}$, $\{51,204\}$ and $\{255\}$. Among the classes of statuses (Z), (S), (A) and (O), those in which $\Phi_f$ has only periodic points, so that the theorem is needed, are the $21$ open classes with finite $\Phi_f$ (Section~\ref{cl:sec:holdouts}) and the classes of $45$, $135$, $63$, $187$, $189$, $60$, $153$, $15$ and $105$.

\begin{remark}[cellwise symmetries]\label{cl:rem:kappa}
Let $f$ be none of $0$, $51$, $204$, $255$. By Theorem~\ref{cl:thm:A} the group $\langle\sh,F_f\rangle\cong\Z^2$ is torsion-free, so a permutation $\pi\ne\mathrm{id}$ of the four cells that commutes with $F_f$ when applied in every cell does not lie in it, and $\langle\sh,F_f,\pi\rangle=\langle\sh,F_f\rangle\times\langle\pi\rangle$. A check over the $23$ nontrivial permutations of $\A$ (\texttt{code/census/classes88.py}) gives the complete list of such symmetries. The complement $\kappa$ is the only one for the ten classes of $24$, $36$, $60$, $90$, $126$, $129$, $153$, $165$, $189$ and $219$, and none of the other classes apart from those of $0$, $51$ and $255$ has any. These ten classes and those of $0$ and $255$ are the self-conjugate classes of \cite[(A.4)]{Takesue1989}; what the check adds is that no other cell permutation commutes. The six pairs listed there consist of a self-conjugate class and the class of its complement, and the correspondence between their orbits \cite[(A.1)--(A.3)]{Takesue1989} implies $F_{1\oplus f}^4=F_f^4$ for a self-conjugate rule $f$; for the pair of $24$ and $189$ this is the relation $F_{189}^4=F_{66}^4$ of Table~\ref{cl:tab:decided}. For $126$ and $219$ the resulting inclusion $\Aut(X_{F_f})\supseteq\Z^2\times\Z/2$ is \cite[Proposition~2.5]{PaperC}; for $24$ and $189$, and for $60$ and $153$, equality holds, by \cite[Theorem~1.1(b)]{PaperA} and \eqref{imp:affine}, and these are the four classes of status (S). The other six of the ten classes have status (G): those of $36$ and $129$ have fixed walls, those of $126$ and $219$ rotating walls, with automorphisms of radius $2$ outside $\langle\sh,F_f\rangle\times\langle\kappa\rangle$ \cite{PaperC}, and those of $90$ and $165$ are settled by \eqref{imp:90}. For $0$, $F_0$ is the transposition $(1\,2)$ of the cells, and the commuting cell permutations are $F_0$, $(0\,3)$ and $\kappa=(0\,3)(1\,2)$; for $51$, $F_{51}$ is the $3$-cycle $(0\,1\,2)$, and the commuting cell permutations are $F_{51}^{\pm1}$; for $255$ they are $F_{255}$, $F_{255}^2=\kappa$ and $F_{255}^3$.
\end{remark}

\subsection{The tables and the proof of Theorem~\ref{thm:main}}\label{cl:sec:main}

The classes are listed with their statuses in Tables~\ref{cl:tab:finite}, \ref{cl:tab:decided} and~\ref{cl:tab:open}. The counts of Theorem~\ref{thm:main} are composed as follows. In status (G), $31$ classes ($81$ rules) have a word of length at most $3$ fixed in every context; the class of $133$ ($2$ rules) has one of length $4$; $5$ classes ($4+4+2+2+2$ rules), those of $59$, $27$, $127$, $179$ and $223$, have hidden-bit or time-periodic walls; the classes of $126$ and $219$ ($1+1$ rules) have rotating walls (in these $39$ classes the walls are those of Takesue \cite[Section~4]{Takesue1989} or are built from them, and the embeddings are those of \cite{PaperC}); the affine classes of $90$ and $165$ ($1+1$ rules) have no wall of any of these kinds and get $(\bigoplus_mS_m)^2$ from marker permutations; the affine class of $255$ ($1$ rule) gets $\bigoplus_mS_m$ from marker permutations of a difference sequence; and the class of $95$ ($2$ rules), which is not affine, gets $\bigoplus_mS_m$ from marker permutations of its light-speed particles. This gives $43$ classes and $104$ rules. In status (Z), the $4$ outer-permutive classes have $16$ rules, the classes of $63$ and $187$ have $4+4$, those of $110$, $14$, $46$ and $79$ have $4+4+2+4$, and the class of $94$ has $2$: $11$ classes and $40$ rules. The $8$ rules of status (S) are $24$, $66$, $189$, $231$, $60$, $102$, $153$ and $195$, and the $6$ rules of status (A) are $15$, $85$, $170$, $240$, $105$ and $150$.

\begin{proof}[Proof of Theorem~\ref{thm:main}]
The classes under reflection and the cellwise complement, the affine and the outer-permutive rules, and the words of length at most $3$ fixed in every context are computed by \texttt{code/census/classes88.py}. (G) For the $31$ classes with a fixed word of length at most $3$, and for the class of $133$, \eqref{imp:walls} and \eqref{imp:133}; for $59$, $27$, $127$, $179$ and $223$, \eqref{imp:hidden}; for $219$ and $126$, \eqref{imp:rot}; for $90$ and $165$, \eqref{imp:90}; for $255$, \eqref{imp:affine}; for $95$, \eqref{imp:95}. (Z) \eqref{imp:perm} for the classes of $30$, $45$, $75$ and $135$; \eqref{imp:half} for the classes of $63$ and $187$; \eqref{imp:110} for the classes of $110$, $14$, $46$ and $79$; \eqref{imp:94} for the class of $94$. (S) \eqref{imp:110} for $24$ and $189$; \eqref{imp:affine} for $60$ and $153$. (A) \eqref{imp:affine} for $15$ and $105$. The remaining $28$ classes are those of Table~\ref{cl:tab:open}; for them the statement on automorphisms of radius at most $3$ is Theorem~\ref{cl:thm:census}, and $\langle\sh,F_f\rangle\cong\Z^2$ is Theorem~\ref{cl:thm:A}, since none of them contains $0$, $51$, $204$ or $255$. The descriptions of the groups in the last sentence of the theorem are those of \eqref{imp:110}, \eqref{imp:affine} and \eqref{imp:95}.

Three remarks on the assembly. The affine classes of $0$ and $51$ have fixed walls and are listed under (G). The affine classes of $90$ and $165$ also have $\kappa$ as their only nontrivial cellwise symmetry (Remark~\ref{cl:rem:kappa}); they are listed under (G) by \eqref{imp:90}. No class of status (Z) is affine; none can have a wall of the kinds of \cite{PaperC}, since such a wall gives status (G).
\end{proof}

\begin{table}[!htbp]
\caption{Status (G): $\bigoplus_mS_m$ embeds in $\Aut(X_{F_f})$. The fixed walls are the walls (i)--(ix) of Takesue \cite[Table~III]{Takesue1989}, and the other walls are built from his periodic walls. For a fixed wall, $W$ is a shortest word fixed by $F_f$ in every context (cells written $0$--$3$); for hidden-bit walls, $k$ is the length of the wall words; for time-periodic walls, each wall word is fixed by $F^P$ in every context; rotating walls are runs of cells that $F$ maps by the $4$-cycle $\pi$ in every context; for $90$ and $165$ the certificate is the marker construction of \eqref{imp:90}, and for $255$ that of \eqref{imp:affine}; these three classes are also affine; for $95$ the certificate is the marker construction of \cite[Lemma~6.8]{PaperE} on the right-moving particles of the two sublattices. The last column gives the result that settles the entry; all but the entry for $95$, which is \cite[Corollary~6.10]{PaperE}, are proved in \cite{PaperC}. The representative is the smallest member.}
\label{cl:tab:finite}
\centering\small\renewcommand{\arraystretch}{0.9}
\begin{tabular}{rlll}
\toprule
rep. & members & certificate & theorem\\
\midrule
0 & 0 & fixed wall $W=0$ & \cite{PaperC}, \eqref{imp:walls}\\
1 & 1, 128 & fixed wall $W=3$ & \cite{PaperC}, \eqref{imp:walls}\\
2 & 2, 8, 16, 64 & fixed wall $W=3$ & \cite{PaperC}, \eqref{imp:walls}\\
3 & 3, 17, 136, 192 & fixed wall $W=3$ & \cite{PaperC}, \eqref{imp:walls}\\
4 & 4, 32 & fixed wall $W=0$ & \cite{PaperC}, \eqref{imp:walls}\\
5 & 5, 160 & fixed wall $W=33$ & \cite{PaperC}, \eqref{imp:walls}\\
6 & 6, 20, 40, 96 & fixed wall $W=33$ & \cite{PaperC}, \eqref{imp:walls}\\
7 & 7, 21, 168, 224 & fixed wall $W=33$ & \cite{PaperC}, \eqref{imp:walls}\\
9 & 9, 65, 130, 144 & fixed wall $W=30$ & \cite{PaperC}, \eqref{imp:walls}\\
10 & 10, 80 & fixed wall $W=30$ & \cite{PaperC}, \eqref{imp:walls}\\
11 & 11, 81, 138, 208 & fixed wall $W=30$ & \cite{PaperC}, \eqref{imp:walls}\\
12 & 12, 34, 48, 68 & fixed wall $W=0$ & \cite{PaperC}, \eqref{imp:walls}\\
18 & 18, 72 & fixed wall $W=3$ & \cite{PaperC}, \eqref{imp:walls}\\
19 & 19, 200 & fixed wall $W=3$ & \cite{PaperC}, \eqref{imp:walls}\\
22 & 22, 104 & fixed wall $W=33$ & \cite{PaperC}, \eqref{imp:walls}\\
23 & 23, 232 & fixed wall $W=33$ & \cite{PaperC}, \eqref{imp:walls}\\
27 & 27, 83, 202, 216 & hidden-bit walls, $k=6$ & \cite{PaperC}, \eqref{imp:hidden}\\
33 & 33, 132 & fixed wall $W=3$ & \cite{PaperC}, \eqref{imp:walls}\\
35 & 35, 49, 140, 196 & fixed wall $W=3$ & \cite{PaperC}, \eqref{imp:walls}\\
36 & 36 & fixed wall $W=00$ & \cite{PaperC}, \eqref{imp:walls}\\
37 & 37, 164 & fixed wall $W=33$ & \cite{PaperC}, \eqref{imp:walls}\\
38 & 38, 44, 52, 100 & fixed wall $W=33$ & \cite{PaperC}, \eqref{imp:walls}\\
39 & 39, 53, 172, 228 & fixed wall $W=33$ & \cite{PaperC}, \eqref{imp:walls}\\
50 & 50, 76 & fixed wall $W=3$ & \cite{PaperC}, \eqref{imp:walls}\\
51 & 51, 204 & fixed wall $W=3$ & \cite{PaperC}, \eqref{imp:walls}\\
54 & 54, 108 & fixed wall $W=33$ & \cite{PaperC}, \eqref{imp:walls}, \eqref{imp:54}\\
55 & 55, 236 & fixed wall $W=33$ & \cite{PaperC}, \eqref{imp:walls}\\
59 & 59, 115, 206, 220 & hidden-bit walls, $k=4$ & \cite{PaperC}, \eqref{imp:hidden}\\
73 & 73, 146 & fixed wall $W=030$ & \cite{PaperC}, \eqref{imp:walls}\\
90 & 90 & affine; markers, $(\bigoplus_mS_m)^2$ & \cite{PaperC}, \eqref{imp:90}\\
95 & 95, 250 & light-speed particles; markers on one family & \cite{PaperE}, \eqref{imp:95}\\
126 & 126 & rotating walls, $c_{i+1}=\kappa(c_i)$ & \cite{PaperC}, \eqref{imp:rot}\\
127 & 127, 254 & time-periodic walls, $P=8$ & \cite{PaperC}, \eqref{imp:hidden}\\
129 & 129 & fixed wall $W=03$ & \cite{PaperC}, \eqref{imp:walls}\\
131 & 131, 137, 145, 193 & fixed wall $W=30$ & \cite{PaperC}, \eqref{imp:walls}\\
133 & 133, 161 & fixed wall $W=0330$ ($k=4$) & \cite{PaperC}, \eqref{imp:133}\\
139 & 139, 209 & fixed wall $W=30$ & \cite{PaperC}, \eqref{imp:walls}\\
147 & 147, 201 & fixed wall $W=303$ & \cite{PaperC}, \eqref{imp:walls}\\
165 & 165 & affine; markers, $(\bigoplus_mS_m)^2$ & \cite{PaperC}, \eqref{imp:90}\\
179 & 179, 205 & time-periodic walls, $P=6$ & \cite{PaperC}, \eqref{imp:hidden}\\
219 & 219 & rotating walls, $c_{i+1}=c_i$ & \cite{PaperC}, \eqref{imp:rot}\\
223 & 223, 251 & time-periodic walls, $P=8$ & \cite{PaperC}, \eqref{imp:hidden}\\
255 & 255 & affine; markers on $x_{i+1}-x_i\in\Z/4$ & \cite{PaperC}, \eqref{imp:affine}\\
\bottomrule
\end{tabular}
\end{table}

\begin{table}[!htbp]
\caption{Statuses (Z), (S) and (A). In the certificate column, ``$\Phi_f$'' means that $\Phi_f$ has a non-periodic point (Section~\ref{cl:sec:Z2}); $\langle\sh,F_f\rangle\cong\Z^2$ holds for every class in the table by Theorem~\ref{cl:thm:A}. The last column gives the result that settles the entry.}
\label{cl:tab:decided}
\centering\footnotesize\renewcommand{\arraystretch}{0.9}
\begin{tabular}{rllll}
\toprule
rep. & members & status & certificate & theorem\\
\midrule
30 & 30, 86, 106, 120 & (Z) & outer-permutive, not affine; $\Phi_f$ & \cite{PaperB}, \eqref{imp:perm}\\
45 & 45, 101, 166, 180 & (Z) & outer-permutive, not affine & \cite{PaperB}, \eqref{imp:perm}\\
75 & 75, 89, 154, 210 & (Z) & outer-permutive, not affine; $\Phi_f$ & \cite{PaperB}, \eqref{imp:perm}\\
135 & 135, 149, 169, 225 & (Z) & outer-permutive, not affine & \cite{PaperB}, \eqref{imp:perm}\\
63 & 63, 119, 238, 252 & (Z) & two neighbours; stationary difference & \cite{PaperB}, \eqref{imp:half}\\
187 & 187, 207, 221, 243 & (Z) & two neighbours; stationary difference & \cite{PaperB}, \eqref{imp:half}\\
62 (class of $110$) & 62, 110, 118, 124 & (Z) & kinks; edge moves; $\Phi_f$ & \cite{PaperA}, \eqref{imp:110}\\
14 & 14, 42, 84, 112 & (Z) & kinks, static kinks; edge moves; $\Phi_f$ & \cite{PaperA}, \eqref{imp:110}\\
46 & 46, 116 & (Z) & kinks; bond bits; $\Phi_f$ & \cite{PaperA}, \eqref{imp:110}\\
79 & 79, 93, 186, 242 & (Z) & kinks of $\sh^{\pm1}F^3$; edge moves; $\Phi_f$ & \cite{PaperA}, \eqref{imp:110}\\
94 & 94, 122 & (Z) & kinks; scattering; Lemma~L; $\Phi_f$ & \cite{PaperE}, \eqref{imp:94}\\
\midrule
24 & 24, 66 & (S) & $\kappa F=F\kappa$; factor onto $46$; $\Phi_f$ & \cite{PaperA}, \eqref{imp:110}\\
189 & 189, 231 & (S) & $\kappa F=F\kappa$; factor onto $46$; $F^4=F_{66}^4$ & \cite{PaperA}, \eqref{imp:110}\\
60 & 60, 102 & (S) & affine; $\Aut=\langle\sh,F\rangle\times\langle\kappa\rangle$ & \cite{PaperC}, \eqref{imp:affine}\\
153 & 153, 195 & (S) & affine; $\Aut=\langle\sh,F\rangle\times\langle\kappa\rangle$ & \cite{PaperC}, \eqref{imp:affine}\\
\midrule
15 & 15, 85, 170, 240 & (A) & affine; $\Aut\cong\Z^2$, index $2$ over $\langle\sh,F\rangle$ & \cite{PaperC}, \eqref{imp:affine}\\
105 & 105, 150 & (A) & affine; $\Aut\cong\Z^3$, with $\mathrm{id}+\sh F$ for $150$ & \cite{PaperC}, \eqref{imp:affine}\\
\bottomrule
\end{tabular}

\medskip
\caption{Status (O), the open classes: $28$ classes, $98$ rules. For each, every automorphism of $X_{F_f}$ of radius at most $3$ is of the form $\sh^aF^b$ (Theorem~\ref{cl:thm:census}). The column $\Phi_f$ records whether $\Phi_f$ has a non-periodic point (Section~\ref{cl:sec:Z2}); for all $28$ classes $\langle\sh,F_f\rangle\cong\Z^2$ holds by Theorem~\ref{cl:thm:A}.}
\label{cl:tab:open}
\centering\footnotesize\renewcommand{\arraystretch}{0.9}
\begin{tabular}{rlcl}
\toprule
rep. & members & $\Phi_f$ & remark\\
\midrule
13 & 13, 69, 162, 176 & -- & \\
25 & 25, 67, 152, 194 & yes & \\
26 & 26, 74, 82, 88 & yes & \\
28 & 28, 56, 70, 98 & -- & \\
29 & 29, 71, 184, 226 & -- & \\
31 & 31, 87, 234, 248 & -- & \\
41 & 41, 97, 134, 148 & -- & \\
43 & 43, 113, 142, 212 & -- & \\
47 & 47, 117, 174, 244 & -- & \\
57 & 57, 99, 156, 198 & -- & permutive in the centre variable\\
58 & 58, 78, 92, 114 & yes & \\
61 & 61, 103, 188, 230 & -- & \\
77 & 77, 178 & yes & \\
91 & 91, 218 & yes & \\
107 & 107, 121, 158, 214 & -- & \\
109 & 109, 182 & -- & \\
111 & 111, 125, 190, 246 & -- & \\
123 & 123, 222 & yes & \\
141 & 141, 163, 177, 197 & -- & \\
143 & 143, 171, 213, 241 & -- & \\
151 & 151, 233 & -- & \\
155 & 155, 203, 211, 217 & yes & \\
157 & 157, 185, 199, 227 & -- & \\
159 & 159, 215, 235, 249 & -- & \\
167 & 167, 173, 181, 229 & -- & \\
175 & 175, 245 & -- & \\
183 & 183, 237 & -- & \\
191 & 191, 239, 247, 253 & -- & \\
\bottomrule
\end{tabular}
\end{table}

\subsection{The census}\label{cl:sec:census}

\begin{theorem}[computer-assisted: SAT enumeration, Section~\ref{sec:method}; one or two rules per class, the other members by Lemma~\ref{lem:classes}]\label{cl:thm:census}
Let $f$ belong to one of the $28$ classes of Table~\ref{cl:tab:open} or to the classes of $14$, $46$, $79$, $94$, $95$ and $110$, and let $R\le3$. Then every invertible block map of radius at most $R$ that commutes with $F_f$ equals $\sh^aF_f^b$ for some $(a,b)$. More precisely, the block maps of radius at most $R$ commuting with $F_f$ are
\begin{enumerate}
\item the $2R^2+2R+1$ maps $\sh^aF_f^b$ with $|a|+|b|\le R$, which are pairwise distinct;
\item the constant maps onto the constant configurations fixed by $F_f$, namely $(0,0)^{\Z}$ if $f(0,0,0)=0$ and $(1,1)^{\Z}$ if $f(1,1,1)=0$;
\item for the class of $95$ and $R=3$ only, $20$ further maps, none of them injective.
\end{enumerate}
\end{theorem}

The proof is a finite computation by the method of Section~\ref{sec:method}; the rules on which it was run and the logs are described in Appendix~\ref{app:census}. For the class of $110$ the same holds for $R=4$ (Theorem~\ref{thm:110}), and for the classes of $110$, $14$, $46$, $79$ and $94$ the first statement holds for every $R$, by \eqref{imp:110} and \eqref{imp:94}. The $20$ further maps for $95$ are identified in \cite[Proposition~6.21]{PaperE}: they are induced by block maps of the golden mean shift that conserve the number of particles and are not injective.

The census was also run for classes whose status is known. For the class of $133$ the commuting block maps of radius at most $3$ are the maps $\sh^aF^b$ with $|a|+|b|\le3$ and nothing else (no constant configuration is fixed, since $f(0,0,0)=f(1,1,1)=1$); so the automorphisms of radius $4$ of \eqref{imp:133} are the first outside $\langle\sh,F\rangle$. For the classes of $24$ and $189$ the census at radius $3$ finds exactly $50$ invertible maps, the $25$ maps $\sh^aF^b$ and their products with $\kappa$, in accordance with \eqref{imp:110}. For the classes of $63$ and $187$, of status (Z) by \eqref{imp:half}, it finds $(2R+1)^2$ invertible maps instead of $2R^2+2R+1$. The reason is that $f_{63}(a,b,c)=\neg(a\wedge b)$ does not read $c$, so $\sh^aF_{63}^b$ depends only on the cells in $[a-|b|,a]$ and can have radius smaller than $|a|+|b|$; for instance $\sh F_{63}^2$ has radius $1$. The pairs $(a,b)$ with $[a-|b|,a]\subseteq[-R,R]$ number $(2R+1)^2$, and an exact comparison of tables shows that the $49$ invertible maps found at $R=3$ are these powers; the same holds for $187$ with $[a,a+|b|]$. So the census agrees with $\Aut=\langle\sh,F\rangle$ for these classes.

Theorem~\ref{cl:thm:census} concerns radius at most $3$. A census at radius $4$ for the open classes is reported in Section~\ref{cl:sec:holdouts}; radius $5$ is out of reach of this encoding, which would need about $10^9$ clauses. The class of $95$ shows what the open status leaves room for: its censuses at radius $3$ (above) and $4$ \cite[Lemma~6.20]{PaperE} contain no automorphism outside $\langle\sh,F\rangle$, and its group contains every finite group, the smallest radius of an automorphism outside $\langle\sh,F\rangle$ being $7$ \eqref{imp:95}.

\subsection{No walls in 110R}\label{cl:sec:nowalls}

The automorphisms of status (G) that come from walls \eqref{imp:walls}--\eqref{imp:rot} all use a cell next to the permuted information whose trajectory does not depend on the context. We recall the kinds of walls of \cite{PaperC}; the fixed walls, and the periodic pairs of cells from which the other kinds are built, are due to Takesue \cite[Section~4]{Takesue1989}. Let $F$ be a reversible cellular automaton of radius $1$ on $\A^{\Z}$.
\begin{itemize}
\item A \emph{fixed wall} is a word $W$ such that $F(c)$ contains $W$ at every position at which $c$ does.
\item A \emph{time-periodic wall} is a family of distinct words $w_0,\dots,w_{P-1}\in\A^k$, indices modulo $P$, such that $c_{[0,k)}=w_s$ implies $F(c)_{[0,k)}=w_{s+1}$, for every $c$ and $s$.
\item A \emph{rotating wall} is a run of cells $c_i$ on which $F$ acts by a fixed permutation $\pi$ of $\A$ in every context, $F(c)_i=\pi(c_i)$, the run being preserved by $F$.
\item A \emph{hidden-bit wall} is given by $k\ge3$, a set $W\subseteq\A^k$, a position $j$ with $1\le j\le k-2$, and an involution $\iota\colon W\to W$ such that $\iota(w)$ differs from $w$ exactly at position $j$. Put $\Occ(c)=\{i:\ c_{[i,i+k)}\in W\}$, and for $i\in\Occ(c)$ let $\iota_ic$ be $c$ with $c_{[i,i+k)}$ replaced by $\iota(c_{[i,i+k)})$. The two basic hypotheses are \textup{(W1)}: if $0\in\Occ(c)$ then $0\in\Occ(F(c))$; and \textup{(H1)}: if $0\in\Occ(c)$ then $F(\iota_0c)=\iota_0F(c)$.
\end{itemize}
The constructions of \cite{PaperC} assume these and further hypotheses; we exclude already the weaker structures defined here. Rule $110$R has none of them, of any length.

\begin{proposition}[no walls in $110$R]\label{cl:prop:nowalls}
Let $F=F_{110}$.
\begin{enumerate}
\item For every word $w$ on an interval $[l,r]$ and every $u\in\A^{(-\infty,l)}$ there are $v,v'\in\A^{(r,\infty)}$ and $t\in\{1,2,3\}$ with $F^t(uwv)_r\ne F^t(uwv')_r$.
\item If $c$ and $c'$ differ in exactly one cell $j$, then $F^t(c)$ and $F^t(c')$ differ in some cell other than $j$ for some $t$ with $1\le t\le6$.
\item Consequently $F$ has no word whose right end cell has a trajectory independent of the context; in particular no word fixed in every context, no time-periodic wall, and no rotating wall. Moreover, for no $k\ge3$ is there a nonempty family $W\subseteq\A^k$ with an involution $\iota$ as above satisfying \textup{(W1)} and \textup{(H1)}. So $110$R has no fixed, time-periodic, rotating or hidden-bit wall of any length.
\end{enumerate}
\end{proposition}

\begin{proof}
Throughout, $f=f_{110}$, and we use three facts read off its table: $f(a,b,\cdot)$ is constant only for $(a,b)=(0,1)$, where it is $1$; $f(1,b,c)=b\oplus c$; and $f(\cdot,0,1)=f(\cdot,1,0)=1$.

(1) Write $Y$ for the diagram of $uwv$. Its values in rows $-1$ and $0$ at columns $i\le r$ do not depend on $v$. Suppose that $F^t(uwv)_r$, which is $(Y(t-1,r),Y(t,r))$, is the same for all $v$, for $t=1,2,3$. For $t=1$: $Y(1,r)=Y(-1,r)\oplus f(Y(0,r-1),Y(0,r),Y(0,r+1))$, and $Y(0,r+1)$ is the free bit $y(v_{r+1})$, so $(Y(0,r-1),Y(0,r))=(0,1)$ and $Y(1,r)=1\oplus Y(-1,r)$. For $t=2$: complementing $x(v_{r+1})=Y(-1,r+1)$ complements $Y(1,r+1)$ and changes no value in row $0$ and no $Y(1,i)$ with $i\le r$. Since $Y(2,r)=Y(0,r)\oplus f(Y(1,r-1),Y(1,r),Y(1,r+1))$, we need $(Y(1,r-1),Y(1,r))=(0,1)$. Then $Y(-1,r)=0$ and $Y(2,r)=1\oplus f(0,1,\cdot)=0$. For $t=3$: complementing $x(v_{r+2})=Y(-1,r+2)$ complements $Y(1,r+2)$ and changes no other value in rows $0$ and $1$, hence no $Y(2,i)$ with $i\le r$; and it complements $Y(2,r+1)=Y(0,r+1)\oplus f(Y(1,r),Y(1,r+1),Y(1,r+2))$, because $Y(1,r)=1$. So $Y(3,r)=Y(1,r)\oplus f(Y(2,r-1),Y(2,r),Y(2,r+1))$ is unchanged only if $(Y(2,r-1),Y(2,r))=(0,1)$, which contradicts $Y(2,r)=0$.

(2) Let $Y=Y_c$, $\Delta=Y_c\oplus Y_{c'}$ and $s_t=\Delta(t,j)$. Suppose, for a contradiction, that $\Delta(t,i)=0$ for all $i\ne j$ and $-1\le t\le T$, where $T=6$.
\begin{enumerate}
\item[(a)] $(s_{t-1},s_t)\ne(0,0)$ for $0\le t\le T$, since otherwise $F^t(c)=F^t(c')$ and $c=c'$.
\item[(b)] If $0\le\tau\le T-1$ and $s_\tau=1$, then $(Y(\tau,j-2),Y(\tau,j-1))=(0,1)$. Indeed, in the relation at $(\tau,j-1)$ the two diagrams differ only in the input $Y(\tau,j)$, and $\Delta(\tau\pm1,j-1)=0$; so $f(Y(\tau,j-2),Y(\tau,j-1),\cdot)$ is constant.
\item[(c)] $s_{\tau+1}=s_{\tau-1}\oplus s_\tau$ for $0\le\tau\le T-1$. If $s_\tau=0$, the relation at $(\tau,j)$ has the same inputs in both diagrams except $Y(\tau-1,j)$. If $s_\tau=1$, then $Y(\tau,j-1)=1$ by (b), and since $f(1,b,c)=b\oplus c$ the relation at $(\tau,j)$ gives $s_{\tau+1}\oplus s_{\tau-1}=1$.
\item[(d)] There is no $\tau$ with $2\le\tau\le T-2$ and $s_\tau=s_{\tau+1}=1$. Suppose there were. By (b), $Y(\tau,j-2)=Y(\tau+1,j-2)=0$ and $Y(\tau,j-1)=Y(\tau+1,j-1)=1$. The relation at $(\tau,j-1)$ gives $Y(\tau-1,j-1)=Y(\tau+1,j-1)\oplus f(0,1,Y(\tau,j))=0$, and the relation at $(\tau,j-2)$ gives $Y(\tau-1,j-2)=Y(\tau+1,j-2)\oplus f(Y(\tau,j-3),0,1)=1$. By (c) at $\tau$ and at $\tau-1$, $s_{\tau-1}=0$ and $s_{\tau-2}=1$, so (b) at $\tau-2\ge0$ gives $Y(\tau-2,j-2)=0$. Then the relation at $(\tau-1,j-2)$ gives $Y(\tau,j-2)=Y(\tau-2,j-2)\oplus f(Y(\tau-1,j-3),1,0)=1$, a contradiction.
\end{enumerate}
By (a), $(s_{-1},s_0)$ is $(0,1)$, $(1,0)$ or $(1,1)$, and by (c) the sequence $s_{-1},\dots,s_6$ is $0,1,1,0,1,1,0,1$, or $1,0,1,1,0,1,1,0$, or $1,1,0,1,1,0,1,1$. These have $s_\tau=s_{\tau+1}=1$ at $\tau=3$, $4$ and $2$ respectively, contrary to (d).

(3) The first statement is (1). A fixed wall, each word of a time-periodic wall, and a rotating wall all have this property: the trajectory of their right end cell is determined by the word. Next let $W\ne\emptyset$, $j$ and $\iota$ satisfy (W1) and (H1), and let $0\in\Occ(c)$. By (W1), $0\in\Occ(F^t(c))$ for all $t\ge0$, so (H1) gives by induction $F^t(\iota_0c)=\iota_0F^t(c)$ for all $t\ge0$. Since $\iota_0$ changes only the cell $j$, the configurations $\iota_0c\ne c$ differ in exactly one cell and their images under every $F^t$ differ in no other cell, contrary to (2).
\end{proof}

The proposition extends to every length the exclusion of fixed words in \eqref{imp:133} and the search for hidden-bit families in \eqref{imp:hidden}, which covers windows of length at most $6$. For words of at most $6$ cells, that the class of $110$ has no local conservation law, and hence no fixed or time-periodic wall, is part of the examination in \cite[Section~4]{Takesue1989}. It does not exclude words fixed in every context by a power $F^P$ alone, which are not one of the wall mechanisms above; a finite search (Appendix~\ref{app:stated}) finds that for $P\le6$ no word has its right end cell fixed by $F^P$ in every context. The statements for the other rules of the class of $110$ follow by reflection and complementation, with left and right exchanged in (1).

\subsection{The open classes}\label{cl:sec:holdouts}

Table~\ref{cl:tab:open} lists the classes for which none of the statuses (G), (Z), (S), (A) is known. The classes of $110$, $14$, $46$, $79$ and $94$ pass the same censuses and have $\Aut(X_{F_f})=\langle\sh,F_f\rangle$; the class of $95$ passes them as far as automorphisms are concerned and has status (G). So the absence of automorphisms outside $\langle\sh,F\rangle$ at small radius does not predict the group, and $\Aut(X_{F_f})=\langle\sh,F_f\rangle$ for the open classes is a question (Question~\ref{q:open}). We record which of the available tools bear on them.

\emph{Walls.} None of the open classes is among the $40$ classes in which Takesue exhibits walls \cite[Section~4]{Takesue1989}. No open class has a word of any length fixed in every context \eqref{imp:133}, and the search of \eqref{imp:hidden} finds hidden-bit and time-periodic families only in the classes of $27$, $59$, $127$, $179$ and $223$. We note that this search does not see families that are unions of several closures, such as the family of length $5$ for rule $126$ in \eqref{imp:rot}, and that walls fixed only by a power of $F$ with larger windows, or longer hidden-bit walls, are not excluded by it. The census would have seen the small automorphisms that walls produce in the known examples: the relabelling $D$ of $54$R has radius $2$ (Table~\ref{tab:census}), the flips of the hidden-bit walls of $59$R appear at radius $2$ (computer-assisted in \cite{PaperC}), and the flip-all maps of the time-periodic walls and of $27$R have radius $2$ and $3$ \cite{PaperC}. This is evidence that short walls of these kinds are absent from the open classes, not a proof: the wall $0330$ of $133$ has length $4$, and its first automorphism outside $\langle\sh,F\rangle$ has radius $4$, beyond the census.

\begin{sloppypar}
\emph{Exact searches.} Finite searches (Appendix~\ref{app:stated}) exclude, for the classes of Table~\ref{cl:tab:open} and the class of $95$: words of length at most $11$, $9$, $7$, $5$, $3$ that are fixed in every context by $\sh^sF^p$ with $p=1,\dots,5$ and $|s|\le p$, and invariant families of type (W1) of the same lengths, except, for $191$, the level sets of its frozen field; hidden-bit walls with words of length at most $10$; and eigenfields $\vartheta$ with $\vartheta\circ F^p=\sh^{-k}\circ\vartheta$ of radius at most $6-p$, for $p=1,\dots,5$ and $|k|\le p$, other than the frozen field of $191$. Eigenfields are the mechanism behind the non-affine commuting maps of $90$R of radius at most $3$ (computer-assisted in \cite{PaperC}). For $p=1$ and $s=0$ an invariant family of type (W1) is a local conservation law in the sense of \cite[Section~4]{Takesue1989}, and for words of length at most $6$ this part of the search is contained in the examination reported there, by which $44$ of the $88$ classes have a local conservation law on at most four cells, and the other $44$ have none on at most six. In the present classification the $44$ classes without one are the open classes other than that of $191$, the classes of status (Z) other than that of $46$, and the classes of $15$, $60$, $90$, $95$, $105$, $153$ and $165$ (Remark~\ref{cl:rem:eigen}). Takesue conjectures that four cells suffice, without proof \cite[p.~391]{Takesue1989}. What the search adds, for the classes on which it is run, is the lengths from $7$ to $11$, and the cases $p\ge2$ and $s\ne0$. The frozen field of $191$ is the local conservation law of $191$R on three cells recorded in \cite[Table~3]{HT1991}. The same searches recover additive conserved quantities for $26$, $77$, $91$, $123$, $175$ and $191$ among the open classes, and for $46$, $94$ and $95$. For densities on at most three adjacent cells these are the quantities of \cite[Tables~3 and~4]{HT1991} (see also \cite{Takesue1989,Takesue1995,Takesue1997}); the quantity of $175$ has a density on four cells, which is beyond the range of \cite{HT1991}; its existence is recorded in \cite[Table~1]{Takesue1997}, which lists one quantity for $175$R with density on four cells and none on fewer, without writing out the density. None of them is known to give an automorphism in an open class. For $95$ they count the particles of each family on each sublattice \cite{PaperE}, and for $46$ the frozen fields are the bond bits used in \cite{PaperA}, that is, the local conservation law of $46$R of \cite[(4.4)--(4.6)]{Takesue1989}.
\end{sloppypar}

\emph{Radius $4$.} A census at radius $4$ (computer-assisted; Appendix~\ref{app:census}), with the method of Section~\ref{sec:method}, covers one rule of each open class and of the classes of $14$, $24$, $46$, $63$, $79$, $94$, $110$, $126$, $133$, $187$ and $189$, $219$. For the open classes and for $14$, $46$, $79$, $94$ and $110$ the commuting block maps of radius at most $4$ are the $41$ maps $\sh^aF^b$ and the constants. For $24$ and $189$ the further solutions are the maps $\kappa\sh^aF^b$, and for $63$ and $187$ they are powers $\sh^aF^b$ of radius smaller than $|a|+|b|$, as an exact comparison of tables confirms. For $133$ there are $27$ further solutions, which are not classified here; automorphisms of radius $4$ outside $\langle\sh,F\rangle$ exist by \eqref{imp:133}. For $126$ and $219$ the enumeration reached its cap. Injectivity is not tested at this radius. For the class of $95$ the census at radius $4$ is \cite[Lemma~6.20, Proposition~6.21(a)]{PaperE}, where it is carried out with two solvers: there are $262$ commuting maps, $220$ of them neither of the form $\sh^aF^b$ nor constant, and none of these is injective.

\emph{The form of the rule.} No open class is permutive in an outer variable: the non-affine outer-permutive rules are the four classes of \eqref{imp:perm}, and the affine ones are the classes of $15$, $105$, $60$, $153$, $90$ and $165$. The rigidity theorem of \cite{PaperB} for automata of any radius applies to rules that are permutive in an outer variable, so it would apply to an open class only through a recoding that makes it outer-permutive, and we know none. The class of $57$ is the only open class that is permutive in the centre variable; the centre-permutive rules are the classes of $51$, $54$, $57$, $60$, $105$, $147$ and $153$. No open class contains a rule that ignores an outer variable: the $28$ such rules lie in classes of status (G), (Z), (S) or (A), and \eqref{imp:half} covers the non-affine ones without fixed walls. The rules that do not read their middle argument are the $16$ rules of the classes of $0$, $5$, $10$, $15$, $90$, $95$, $165$, $175$ and $255$ \cite[p.~386]{Takesue1989}; the class of $175$ is the only open one, and for it see \cite[Remark~6.22]{PaperE}.

\emph{Light-speed fixed points.} The proofs of status (Z) in \cite{PaperA,PaperE} begin with a homomorphism $\Aut(X_{F_f})\to\Z^2$ that extends $\sh^aF^b\mapsto(a,b)$, read off from the action of an automorphism on kinks among the fixed points of $\sh^{\pm1}F_f$, that is, on non-periodic points of $\Phi_f$. For $7$ of the open classes $\Phi_f$ has non-periodic points (Table~\ref{cl:tab:open}): for the class of $25$ they come from a path between two cycles, as for $110$, and for the other six $\Phi_f$ has positive entropy, as for $94$ and $95$. For the remaining $21$ open classes $\Phi_f$ is finite, and for $9$ of them ($141$, $143$, $151$, $157$, $159$, $167$, $175$, $183$, $191$) it is empty: there $f(0,0,0)=f(1,1,1)=1$ and no constant configuration is fixed by $F$. For these $21$ classes no such homomorphism is known.

\emph{Blocking words.} In the wall constructions of \cite{PaperC} the automorphisms outside $\langle\sh,F\rangle$ come from a blocking word (Section~\ref{sec:blocking}) with an autonomous degree of freedom next to it. The absence of blocking words does not imply that $\Aut(X_{F_f})$ is virtually $\Z^2$: not for affine rules, by \eqref{imp:90}, and not for rules that are not affine, by \eqref{imp:95}. It is not known which open classes have blocking words.

\emph{Window graphs and eigenfields.} By Lemma~\ref{cl:lem:eigen}, the condition $(\Lambda^{\pm})$ of Section~\ref{sec:blocking} says that $F$ has no eigenfield of speed $\pm1$, which the eigen-factors of $90$R violate \cite{PaperC}, and the zero-speed condition says that $F$ has no frozen field. For $95$R the graphs $W^{(k,p)}_M$ are strongly connected for all $(k,p)$ and all $M$ \eqref{imp:95}.

\begin{remark}[window graphs of the classes; computer-assisted]\label{cl:rem:eigen}
A computation of the graphs $W^{(\pm1,1)}_M$ and $W^{(0,1)}_M$ for $M\le6$ (Appendix~\ref{app:stated}) gives the following. Every class of status (G) except those of $90$, $165$ and $95$ fails the zero-speed condition at some $M\le3$; $90$ and $165$ satisfy it for $M\le6$ and fail $(\Lambda^{\pm})$. The classes of status (Z) other than that of $46$, and $27$ of the $28$ open classes, satisfy both conditions for $M\le6$; the exceptions $46$ and $191$ fail the zero-speed condition and have frozen fields. In status (S), the classes of $24$ and $189$ fail the zero-speed condition, at $M=1$ and $M=2$, and the affine classes of $60$ and $153$ satisfy both conditions for $M\le6$, as do the classes of $15$ and $105$ of status (A). The class of $95$ satisfies both conditions for every $M$, by \eqref{imp:95}. For the zero-speed graphs with $M\le5$, that is, for words of at most $6$ cells, these statements reproduce the examination of local conservation laws in \cite[Section~4]{Takesue1989}; the graphs of speed $\pm1$ are not considered there. These computations are not used elsewhere.
\end{remark}

\subsection{Reach of the edge-move method}\label{cl:sec:reach}

The proof of \eqref{imp:110} in \cite{PaperA} works at the right edge of configurations that agree with a spatially periodic background on a right half-line. We record where this method stops on the open classes. Let $f$ be an elementary rule and $F=F_f$. A configuration $B$ (in the applications, a spatially periodic one) has a \emph{stuck right edge at $r$} if for every $t\in\Z$ and every configuration $c$ that agrees with $F^tB$ on $(r,\infty)$ and differs from it at $r$, the configurations $F(c)$ and $F^{-1}(c)$ agree with $F^{t+1}B$ and $F^{t-1}B$ on $(r,\infty)$ and differ from them at $r$. Stuck left edges are defined in the same way with $(-\infty,r)$, and the reflection $\mathrm{rev}$ exchanges the two. For $110$, $42$ and $46$ the constant $(0,0)^{\Z}$ has a stuck right edge at every $r$, and for $93$ a background of spatial period $2$ has one \cite{PaperA}. Let
\[
N_f=\{(b,c)\in\{0,1\}^2:\ f(0,b,c)=f(1,b,c)\}
\]
be the set of contexts in which $f$ ignores its left input, and for $(b,c)\in N_f$ let $g_f(b,c)$ be the common value.

\begin{lemma}\label{cl:lem:stuckcols}
Let $B$ have a stuck right edge at $r$, and put $u_t=Y_B(t,r+1)$, $v_t=Y_B(t,r+2)$, $w_t=Y_B(t,r+3)$, $z_t=Y_B(t,r+4)$. Then for every $t\in\Z$:
\begin{enumerate}
\item $(u_t,v_t)\in N_f$;
\item $u_{t+1}\oplus u_{t-1}=g_f(u_t,v_t)$;
\item $v_{t+1}\oplus v_{t-1}=f(u_t,v_t,w_t)$ and $w_{t+1}\oplus w_{t-1}=f(v_t,w_t,z_t)$.
\end{enumerate}
\end{lemma}

\begin{proof}
Fix $t$, let $d=F^tB$, and let $c$ be $d$ with the bit $y(d_r)$ complemented. Then $c$ agrees with $d$ on $(r,\infty)$ and differs from it at $r$, so $F(c)_{r+1}=F(d)_{r+1}$. By \eqref{eq:F} the second bits of these cells are $x(d_{r+1})\oplus f(\neg y(d_r),y(d_{r+1}),y(d_{r+2}))$ and $x(d_{r+1})\oplus f(y(d_r),y(d_{r+1}),y(d_{r+2}))$, and the $y$-row of $d$ is row $t$ of $Y_B$; this gives (1). Parts (2) and (3) are the spacetime relation \eqref{eq:rel} of $Y_B$ at $(t,r+1)$, $(t,r+2)$ and $(t,r+3)$, with (1) used in (2).
\end{proof}

\begin{proposition}\label{cl:prop:nostuck}
Let $f$ be a rule in one of the $12$ classes of $77$, $91$, $94$, $95$, $107$, $109$, $123$, $141$, $151$, $167$, $175$ and $183$. Then no configuration has a stuck right edge or a stuck left edge for $F_f$; in particular no spatially periodic one has.
\end{proposition}

\begin{proof}
The complement $\kappa$ commutes with $\sh$ and conjugates $F_f$ to $F_{\kappa f}$, so it maps a configuration with a stuck right edge for $F_f$ to one for $F_{\kappa f}$; and $\mathrm{rev}$ maps stuck left edges for $F_f$ to stuck right edges for $F_{f^r}$, which lies in the same class. So it suffices to exclude stuck right edges for one rule of each pair $\{f,\kappa f\}$ in these classes, namely for
\[
91,\ 94,\ 95,\ 107,\ 121,\ 109,\ 123,\ 151,\ 167,\ 173,\ 175,\ 183,\ 77,\ 141,\ 163.
\]
Suppose that $B$ has a stuck right edge at $r$, and use the notation of Lemma~\ref{cl:lem:stuckcols}. The values of $f$ used below are read off its table.

(a) For the first twelve rules, $g_f\equiv1$ on $N_f$, so by (2) $u_{t+1}\ne u_{t-1}$ for every $t$, and $u$ takes both values. For $91$, $94$, $107$, $121$, $109$, $123$, $151$, $167$, $173$, $183$ the set $N_f$ is $\{00\}$, $\{10\}$, $\{01\}$, $\{00\}$, $\{10\}$, $\{00,01\}$, $\{00\}$, $\{01\}$, $\{11\}$, $\{00,01\}$, whose elements share their first coordinate, so $u$ is constant by (1), a contradiction. For $95$, $N_f=\{00,10\}$, and for $175$, $N_f=\{01,11\}$; in both cases $v$ is constant by (1), so $v_{t+1}\oplus v_{t-1}=0$, and at a time $t$ with $u_t=0$ part (3) requires $f(0,v_t,w_t)=0$, while $f_{95}(0,0,\cdot)=1$ and $f_{175}(0,1,\cdot)=1$.

(b) For $77$ and $163$, $N_f=\{01,10\}$, so $v_t=\neg u_t$ and (2), (3) give $f(u_t,v_t,w_t)=g_f(u_t,v_t)$. At a time with $u_t=0$ this reads $f(0,1,w_t)=g_f(0,1)$, while $f_{77}(0,1,\cdot)=1\ne0=g_{77}(0,1)$ and $f_{163}(0,1,\cdot)=0\ne1=g_{163}(0,1)$. So $(u_t,v_t)=(1,0)$ for all $t$, and (2) requires $g_f(1,0)=0$. For $77$, $g_{77}(1,0)=1$. For $163$, $g_{163}(1,0)=0$, and (3) gives $f_{163}(1,0,w_t)=0$ for all $t$, so $w\equiv0$ since $f_{163}(1,0,1)=1$; then the second relation in (3) reads $0=f_{163}(0,0,z_t)$, while $f_{163}(0,0,\cdot)=1$.

(c) For $141$, $N_f=\{01,11\}$, so $v\equiv1$ and $f_{141}(u_t,1,w_t)=0$ by (3). Since $f_{141}(0,1,\cdot)=1$, $u\equiv1$, and (2) requires $g_{141}(1,1)=0$, while $g_{141}(1,1)=1$.
\end{proof}

In these $12$ classes, of which all but those of $94$ and $95$ are open, the construction of \cite{PaperA} cannot start: there is no edge at which to glue a kink. They are not the only classes without stuck edges: for the affine classes of $90$, $105$ and $165$ the rule is permutive in both outer variables, $N_f$ is empty, and part (1) of Lemma~\ref{cl:lem:stuckcols} excludes stuck edges at once. For the remaining open classes the obstruction is located by computation.

\begin{proposition}[computer-assisted]\label{cl:prop:reach}
Let $f$ be a rule of one of the $11$ classes of $28$, $29$, $41$, $43$, $47$, $57$, $61$, $111$, $143$, $155$ and $157$, let $B$ be a spatially periodic configuration of spatial period at most $6$ with a stuck right edge, and let $1\le b\le4$ and $0<|a|\le b$. If $B\in\mathrm{Fix}(\sh^aF_f^b)$, then $B$ is the only point of $\mathrm{Fix}(\sh^aF_f^b)$ that agrees with $B$ on a left half-line.
\end{proposition}

The proof is an exact finite search (Appendix~\ref{app:reach}): in the graph of the windows of $\mathrm{Fix}(\sh^aF_f^b)$, every window of $B$ has exactly one outgoing edge. In particular, for these rules, periods and speeds there is no kink with such a background on its left, and nothing to glue next to the edge.

\begin{remark}[sampled computations]\label{cl:rem:reach}
The following computations use sampled left contexts (Appendix~\ref{app:reach}); they are observations, not proofs. In \cite{PaperA}, a kink glued next to a stuck edge leaves, after a bounded time, a configuration with the same edge that is changed only near the edge; this defines a \emph{move}, a partial map on the states of the last $k$ cells before the edge, and the proof needs the orbits of the moves to reach the states from which the next column to the left can be read.
\begin{enumerate}
\item For the class of $159$ there are kinks of speed $1/2$ with a background with a stuck edge on their left, but no kink glued next to the edge was found to give a move.
\item For the classes of $13$, $25$, $26$, $58$ and $191$ there are backgrounds with a stuck edge, kinks and moves, but the orbits of the moves on the last $k\le4$ cells do not reach the states from which the next column is read.
\item For the class of $31$ there are backgrounds, kinks of speed $1/3$ and moves, and for every $k\le6$ the moves on the last $k$ cells have two orbits for one value of the cell beyond the edge and a single orbit for each other value, as for $110$. A proof that a block map commuting with $F_{31}$ and fixing the kinks commutes with the moves is missing.
\end{enumerate}
With Propositions~\ref{cl:prop:nostuck} and~\ref{cl:prop:reach} this accounts for the $28$ open classes: $10$, $11$, $1$, $5$ and $1$.
\end{remark}
